% Part: sets-functions-relations
% Chapter: arithmetization
% Section: checking-details
% Hindi translation (hi-Deva-IN) of Open Logic commit 9620cc73f9c8e0ad003c514a5d3748f29611c4c0.
%
\documentclass[../../../include/open-logic-section]{subfiles}


\begin{document}
	
\olfileid[hi]{sfr}{arith}{check}
\olsection{क्रमित वलय और क्षेत्र}

इस पूरे अध्याय में हमने दावा किया कि कुछ परिभाषाएँ ``अपेक्षित ढंग से'' व्यवहार
करती हैं। इस तकनीकी परिशिष्ट में हम स्पष्ट करेंगे कि इसका क्या अर्थ है, और
(रूपरेखा में दिखाएँगे कि) ये परिभाषाएँ सचमुच ``ठीक'' व्यवहार करती हैं।

\olref[int]{sec} में हमने $\Int$ पर योग और गुणन परिभाषित किए थे। हमें दिखाना
है कि ये परिभाषित संक्रियाएँ $\Int$ को वही संरचना देती हैं जो हम उससे
``चाहेंगे''। विशेषतः यह संरचना एक क्रमविनिमेय वलय की है।

\begin{defn}
	कोई \emph{क्रमविनिमेय वलय} ऐसा समुच्चय $S$ है जिसमें विशिष्ट अवयव $0$ और
	$1$, तथा संक्रियाएँ $+$ और $\times$ दी गई हों और जो इन आठ सूत्रों को
	संतुष्ट करे:
	\begin{align*}
		\emph{साहचर्यता}&&a + (b+ c) & = (a + b) + c \\
		&& (a \times b) \times c & = a \times (b\times c)\\
		\emph{क्रमविनिमेयता}&&a + b &= b+ a  \\
		&&  a \times b&= b\times a\\ 
		\emph{तत्समक अवयव}&&a + 0 &= a \\
		&& a \times 1 &= a\\
		\emph{योगात्मक प्रतिलोम}&&(\exists b\in S)0&=a + b\\
		\emph{वितरणशीलता}&&a \times (b+ c ) &= (a \times b) + (a \times c)
	\end{align*}
	यहाँ निहित रूप से इन सभी सूत्रों पर $S$ तक सीमित सार्वत्रिक परिमाणक लगे
	हैं। यह भी ध्यान दें कि यहाँ $0$ और~$1$ नाम वाले अवयवों का समान नाम वाली
	प्राकृत संख्याएँ होना आवश्यक नहीं है।
\end{defn}

अतः यह जाँचने के लिए कि पूर्णांक क्रमविनिमेय वलय बनाते हैं, केवल इन आठ शर्तों
की जाँच करनी है। किसी शर्त को स्थापित करना कठिन नहीं है, किंतु काम कुछ लंबा
है। उदाहरण के लिए, योग के मामले में \emph{साहचर्यता} इस प्रकार सिद्ध होती है:

\begin{proof} 
$i, j, k \in \Int$ नियत करें। तब $a_1, b_1, a_2, b_2, a_3, b_3 \in
\Nat$ ऐसे हैं कि $i = \equivrep{a_1, b_1}{}$ और $j =
\equivrep{a_2,b_2}{}$ और $k = \equivrep{a_3, b_3}{}$। (पठनीयता के लिए हम
``$\equivrep{x, y}{}$'' लिखेंगे, न कि ``$\equivrep{x, y}{\Intequiv}$'';
इस पूरे खंड में ऐसा ही करेंगे।) अब:
\begin{align*}
	i + (j + k) &= \equivrep{a_1, b_1}{}+(\equivrep{a_2, b_2}{} + \equivrep{a_3, b_3}{}) \\
	&= \equivrep{a_1,  b_1}{} + \equivrep{a_2+a_3, b_2+b_3}{}\\
	&= \equivrep{a_1 + (a_2 + a_3), b_1 + (b_2 + b_3)}{}\\
	&= \equivrep{(a_1 + a_2) + a_3, (b_1 + b_2) + b_3}{}\\
	&= \equivrep{a_1 + a_2, b_1 + b_2}{} + \equivrep{a_3, b_3}{}\\
	&= (\equivrep{a_1, b_1}{} + \equivrep{a_2, b_2}{}) + \equivrep{a_3, b_3}{}\\
	&= (i+j) + k
\end{align*}
जहाँ हमने $\Nat$ पर योग के व्यवहार का निर्बाध उपयोग किया है।
\end{proof}

इसी प्रकार \emph{योगात्मक प्रतिलोम} का प्रमाण यह है:

\begin{proof}
$i \in \Int$ नियत करें, ताकि $i = \equivrep{a,b}{}$ हो, जहाँ $a,b \in
\Nat$। $j = \equivrep{b,a}{} \in \Int$
रखें। प्राकृत संख्याओं के व्यवहार का उपयोग करने पर $(a+b) + 0 = 0 + (a+b)$,
इसलिए परिभाषा से $\tuple{a+b, b+a} \sim_\Int \tuple{0,0}$, और फलतः
$\equivrep{a+b, b+a}{} = \equivrep{0, 0}{} = 0_\Int$। अब $i + j =
\equivrep{a,b}{}+\equivrep{b,a}{}=\equivrep{a+b, b+a}{}= \equivrep{0,
0}{} = 0_\Int$।
\end{proof}

और \emph{वितरणशीलता} का प्रमाण यह है:

\begin{proof}
ऊपर की तरह $i = \equivrep{a_1, b_1}{}$ और $j =
\equivrep{a_2,b_2}{}$ और $k = \equivrep{a_3, b_3}{}$ नियत करें। अब:
\begin{align*}
	i \times (j + k) 
	&= \equivrep{a_1, b_1}{} \times (\equivrep{a_2,b_2}{} + \equivrep{a_3, b_3}{})\\
	&= \equivrep{a_1, b_1}{} \times \equivrep{a_2 + a_3,b_2+b_3}{}\\
	&= \equivrep{a_1  (a_2 + a_3) + b_1  (b_2+b_3), a_1  (b_2 + b_3) + b_1 (a_2 + a_3)}{}\\
	&= \equivrep{a_1 a_2 + a_1a_3 + b_1 b_2+b_1b_3, a_1 b_2 + a_1b_3 + a_2b_1 + a_3b_1}{}\\		
%		&= \equivrep{a_1 a_2 + b_1b_2 + a_1a_3 + b_1 b_2+b_1b_3, a_1 b_2 + a_1b_3 + a_2b_1 + a_3b_1}{}\\		
	&= \equivrep{a_1a_2 + b_1b_2, a_1b_2 + a_2b_1}{} + \equivrep{a_1a_3 + b_1b_3, a_1b_3 + a_3b_1}{}\\
	&= (\equivrep{a_1, b_1}{} \times \equivrep{a_2,b_2}{}) + (\equivrep{a_1, b_1}{} \times  \equivrep{a_3, b_3}{})\\
	&= (i \times j) + (i \times k)
\end{align*}
\end{proof}

शेष पाँच शर्तों के प्रमाण को अभ्यास के रूप में छोड़ते हैं। उन्हें सिद्ध करने
पर यह दिख जाएगा कि $\Int$ एक क्रमविनिमेय वलय है, अर्थात योग और गुणन
(जैसा उन्हें परिभाषित किया गया है) अपेक्षित ढंग से व्यवहार करते हैं।

\begin{prob}
सिद्ध कीजिए कि $\Int$ एक क्रमविनिमेय वलय है।
\end{prob}

लेकिन हमारा काम अभी समाप्त नहीं हुआ। $\Int$ पर योग और गुणन के साथ हमने एक
क्रम-संबंध $\leq$ भी परिभाषित किया था, और हमें जाँचना है कि वह भी अपेक्षित
ढंग से व्यवहार करता है। अधिक विस्तार से, दिखाना है कि $\Int$ एक
\emph{क्रमित} वलय है।\footnote{\olref[sfr][rel][ord]{def:linearorder} से याद
	करें कि संपूर्ण क्रम ऐसा संबंध है जो स्वतुल्य, संक्रामक, प्रति-सममित और
	संबद्ध हो। क्रम-संबंधों के संदर्भ में संबद्धता को कभी-कभी
	\emph{त्रिविकल्प} कहते हैं, क्योंकि किन्हीं $a$ और $b$ के लिए $a \leq b \lor a
	= b \lor a \geq b$ होता है।}

\begin{defn}
कोई \emph{क्रमित वलय} ऐसा क्रमविनिमेय वलय है जिस पर एक संपूर्ण क्रम-संबंध
$\leq$ भी दिया गया हो और जो इन शर्तों को संतुष्ट करे:
\begin{align*}
	a \leq b &\lif a + c \leq b + c\\
	(a \leq b \land 0 \leq c) &\lif a \times c \leq b \times c
\end{align*}
\end{defn}

\begin{prob}
सिद्ध कीजिए कि $\Int$ एक क्रमित वलय है।
\end{prob}

पहले की तरह यह दिखाना लंबा, किंतु नियमित काम है कि निर्मित $\Int$ एक क्रमित
वलय है। यह काम आपके लिए छोड़ते हैं।

इससे पूर्णांकों का काम पूरा हुआ। अब परिमेय संख्याओं के बारे में बहुत मिलती-जुलती
बातें दिखानी हैं। विशेषतः योग, गुणन और क्रम की हमारी दी हुई परिभाषाओं के अधीन
परिमेय संख्याओं को एक क्रमित \emph{क्षेत्र} सिद्ध करना है; ये संक्रियाएँ हैं
$+$, $\times$, और $\leq$:
\begin{defn}\ollabel{orderedfield}
कोई \emph{क्रमित क्षेत्र} ऐसा क्रमित वलय है जो यह शर्त भी संतुष्ट करे:
\begin{align*}
	\emph{गुणात्मक प्रतिलोम}& & (\forall a \in S \setminus \{0\})(\exists b \in S) a\times b& = 1
\end{align*}
\end{defn}

एक बार यह दिखा देने पर कि $\Int$ एक क्रमित वलय है, यह दिखाना आसान किंतु लंबा
है कि $\Rat$ एक क्रमित क्षेत्र है।

\begin{prob}
सिद्ध कीजिए कि $\Rat$ एक क्रमित क्षेत्र है।
\end{prob}

पूर्णांकों और परिमेय संख्याओं का काम होने के बाद केवल वास्तविक संख्याएँ बचती
हैं। विशेषतः दिखाना है कि $\Real$ एक \emph{पूर्ण} क्रमित क्षेत्र है, अर्थात
पूर्णता गुणधर्म वाला क्रमित क्षेत्र। \olref[cuts]{realcompleteness} ने स्थापित
किया कि $\Real$ में पूर्णता गुणधर्म है। किंतु अभी यह जाँचने का (लंबा) काम
बाकी है कि $\Real$ एक क्रमित क्षेत्र है।

\emph{उस} लंबे अभ्यास में उतरने से पहले कुछ और ``तत्काल'' बातों की जाँच करनी होगी।
उदाहरण के लिए, यह सुनिश्चित करना है कि परिभाषित $\alpha + \beta$ सचमुच एक
\emph{विच्छेद} है, किन्हीं विच्छेदों $\alpha$ और $\beta$ के लिए। इसका प्रमाण
यह है:

\begin{proof}	
चूँकि $\alpha$ और $\beta$ दोनों विच्छेद हैं, इसलिए $\alpha + \beta = \Setabs{p
+ q}{p \in \alpha \land q \in \beta}$, $\Rat$ का रिक्तेतर उचित उपसमुच्चय
है। अब मान लें कि $x < p + q$, किन्हीं $p \in \alpha$ और $q \in \beta$ के लिए।
तब $x - p < q$, इसलिए $x - p \in \beta$, और $x = p + (x - p) \in
\alpha + \beta$। अतः $\alpha + \beta$, $\Rat$ का आरंभिक खंड है। अंततः
किसी भी $p + q \in \alpha + \beta$ के लिए, चूँकि $\alpha$ और $\beta$
दोनों विच्छेद हैं, ऐसे $p_1 \in \alpha$ और $q_1 \in \beta$ हैं कि
$p < p_1$ और $q < q_1$; इसलिए $p + q < p_1 + q_1 \in \alpha +
\beta$; अतः $\alpha + \beta$ का कोई अधिकतम अवयव नहीं है।
\end{proof}

ऐसे ही प्रयासों से आप जाँच सकते हैं कि $\alpha - \beta$, $\alpha \times \beta$
और $\alpha \div \beta$ विच्छेद हैं (अंतिम मामले में उस स्थिति को छोड़कर जहाँ
$\beta$ शून्य-विच्छेद है)। इसे भी आपके लिए छोड़ते हैं।

\begin{prob}
सिद्ध कीजिए कि $\Real$ एक क्रमित क्षेत्र है।
\end{prob}

एक छोटा खुला सिरा अभी समेटना है। \olref[cuts]{sec} में हमने सुझाव दिया कि
$\sqrt{2} =
\Setabs{p \in \Rat}{p < 0 \text{ अथवा }p^2 < 2}$ लिया जा सकता है। किंतु यह
दिखाना आवश्यक है कि यह समुच्चय एक \emph{विच्छेद} है। इसका प्रमाण यह है:

\begin{proof}
स्पष्टतः यह परिमेय संख्याओं का रिक्तेतर उचित आरंभिक खंड है; अतः केवल यह दिखाना
पर्याप्त है कि इसका कोई अधिकतम अवयव नहीं है। विशेषतः यह दिखाना पर्याप्त है कि
जहाँ $p$ ऐसा धनात्मक परिमेय है जिसके लिए $p^2 < 2$ और $q =
\frac{2p+2}{p+2}$, वहाँ $p < q$ और $q^2 < 2$ दोनों होते हैं। $p < q$ देखने
के लिए बस ध्यान दें:
\begin{align*}
	p^2 &< 2\\
	p^2 + 2p &< 2 + 2p\\
	p(p + 2) &< 2 + 2p\\
	p &< \tfrac{2+2p}{p+2} = q
\end{align*}
$q^2 < 2$ देखने के लिए बस ध्यान दें:
\begin{align*}
	p^2 &< 2\\
	2p^2 + 4p + 2 &< p^2 + 4p+ 4\\
	4p^2 + 8p + 4 &< 2(p^2 + 4p + 4)\\
	(2p+2)^2 & <2(p+2)^2\\
	\tfrac{(2p+2)^2}{(p+2)^2} &< 2\\
	q^2 &< 2
\end{align*}
\end{proof}
\end{document}
